\scanpage{6}
\mspage{\handwritten{4}}

\deleted{pleinement fidèle (resp} une équivalence). Même énoncé si on met dans le
coup un \deleted{Anneau loc} systême local d'anneaux sur X'. Enfin, f induit \deleted{un}
\inserted{\deleted{foncteur \unsure{soit équiv.|plein fid} (resp} une équivalence).}
\deleted{isomorphisme} entre la catégorie des coefficients locaux sur C et celle
des coéff. locaux sur C' \deleted{(?? peut être}
\inserted{\deleted{cette propriété doit}}
\deleted{\unsure{il|l} doit il être \unsure{mise|mis} dans les
hypothèses préliminaires sur f ?)} Principe de démonstration: on commence
par prouver ce dernier résultat, en notant que \deleted{dans le ca} si un topos est
localement hom. trivial, il est loca connexe et loc. simplement connexe,
d'où une bonne théorie du $\handwritten{\pi}_{0}$ et du $\handwritten{\pi}_{1}$ (qui sont ici discrets), et on est
ramené \deleted{pour} à un cas particulier du critère d'homotopisme de Artin et
Mazur, savoir un critère \begin{tabular}[b]{@{}l@{}}\inserted{cohomologique}\\[-.6ex]\begin{tikzpicture}[x=1cm,y=1cm,baseline] \draw (0,0) .. controls (-.17,.15) and (-.12,.43) .. (.13,.64); \draw (.18,.69) .. controls (.55,.95) and (1.08,.96) .. (1.55,.87) .. controls (1.92,.8) and (2.17,.81) .. (2.35,.8); \draw[dotted] (2.43,.8) -- (3.42,.84); \draw (3.51,.85) -- (3.73,.88); \draw[dotted] (3.77,.9) -- (3.94,1.02); \draw (3.99,1.08) .. controls (4.04,1.25) and (4.03,1.37) .. (3.94,1.45); \end{tikzpicture}\end{tabular} pour qu'un homomorhisme de groupes $G\handwritten{\to} H$ (ici de
groupes $\handwritten{\pi}_{1}$ de X,X') soit un isomorphisme: il doit induire des isomorphismes sur les $H^{o}$ et $H^{1}$ (y inclus dans le cas non commutatif...). On prouve
la pleine fidélité en se ramenant par la suite spectrale aux cas des $\mathrm{Ext}^{\unsure{i}}$
de coéfficients locaux. Par la suite spectrale encore, cela résultera du
fait suivant: si X\deleted{'} est \unsure{\deleted{l}h|\deleted{h}l}omalement hom. trivial, alors la catégorie des
faisceaux abéliens loc. constants est stable par $\underline{\mathrm{Ext}}^{i}$ , et le foncteur
$M\handwritten{\mapsto} M_{X}$ \deleted{des fai\unsure{sce}a\unsure{ux} constants} (Ab) dans $X_{ab}$ commute aux dits $\mathrm{Ext}^{i}$ .
En fait, X et X' étant loc.homp\begin{tikzpicture}[x=1cm,y=1cm,baseline=0pt]\draw[line width=.4pt] (0,.35)--(-.08,.85);\draw[line width=.4pt] (.02,-.08)--(.1,-.65);\end{tikzpicture}triviaux, les conditions suivantes sur f
seront équivalentes:

a) f est\unsure{!)|\deleted{'}} un homotopisme, i.e. induit pour tout systême local (pas
néc commutatif) sur X'un isomorphisme sur les $H^{i}$\deleted{pour i \unsure{m}}.

b) \deleted{Même chose ave} f induit une équivalence $D^{b}_{lc}(X')\handwritten{\longrightarrow} D^{b}_{lc}(X)$ .

a\handwritten{'}) f induit une équivalence entre la catégorie des systêmes locaux
abéliens sur X' et X, et des isomorphismes sur les $H^{\unsure{i|1}}$ correspondants (do\unsure{n}
on ne prend ici que des coéfficients commutatifs).

\margin{\handwritten{?}} J'ignore si on peut dans a) se borner aux systêmes locaux commutatif\unsure{s}

\deleted{i e \unsure{dan}}
L'équivalence entre a) et b) fournit une première justification ou
motivation pour définir des types d'homotopie via la catégorie $D\handwritten{^{b}_{lc}(X)}$,

\margin{\handwritten{5}}
